\documentclass[../main.tex]{subfiles}
\graphicspath{{\subfix{../images/}}}

\begin{document}
    \section{Manual de Uso}
    El programa adjunto es una aplicación de línea de comandos desarrollada en
    lenguaje C que permite validar cadenas mediante máquinas de Turing
    previamente definidas. El sistema cuenta con dos máquinas:

    \begin{itemize}
        \item M2: Reconoce el Lenguaje 2
        \item M8: Reconoce el Lenguaje 8
    \end{itemize}

    El usuario debe ingresar la cadena a evaluar y el lenguaje al que se desea
    comprobar si pertenece; el programa seleccionará automáticamente la máquina de
    Turing correspondiente.

    Se envió el ejecutable compilado listo para ejecutarse al hacer doble click.
    Después se abrirá una ventana de consola negra. Ahí ya se puede escribir
    comandos.

    Las opciones para los comandos siguen la convención UNIX: cada opción tiene
    una forma corta (p.ej. \code{-i}) y una forma larga equivalente (p.ej.
    \code{--in}), separadas por coma en los ejemplos de sintaxis. Se debe
    escoger una de las dos. No deben utilizarse ambas a la vez.

    \section{Comandos disponibles}

    \subsection{Comando \texttt{validate}}
    Este comando permite validar cadenas utilizando una máquina de Turing
    específica. Muestra en pantalla ``Cadena válida'' si la cadena dada pertenece al lenguaje especificado o ``Cadena inválida'' en caso contrario. Admite dos modos de uso:

    \textbf{Sintaxis 1 — validación de una cadena:}

    \begin{lstlisting}[style=console]
    validate <cadena> -i, --in <lenguaje>
    \end{lstlisting}

    \begin{description}
        \item[\code{<cadena>}] Cadena a evaluar, por ejemplo: \code{aabbb}.
        \item[\code{<lenguaje>}] Indica el lenguaje al que debe pertenecer la cadena dada.
    \end{description}

    \textbf{Sintaxis 2 — modo prompt:}

    \begin{lstlisting}[style=console]
    validate -p, --prompt <lenguaje>
    \end{lstlisting}

    \begin{description}
        \item[\code{-p}, \code{--prompt}] Inicia un prompt interactivo en el
        que se pueden ingresar cadenas una tras otra para validarlas sobre el
        lenguaje especificado. Utilize \code{exit} para salir de este modo.
        \item[\code{<lenguaje>}] Selecciona la máquina de Turing. Valores
        posibles: \code{m2}, \code{m8}.
    \end{description}

    \textbf{Sintaxis 3 — listado de máquinas:}

    \begin{lstlisting}[style=console]
    validate -l, --list
    \end{lstlisting}

    \begin{description}
        \item[\code{-l}, \code{--list}] Se mostrará las máquinas disponibles en pantalla.
    \end{description}

    \subsection{Comando \texttt{clear}}

    El programa incluye un comando para limpiar la pantalla. Este comando elimina
    el contenido mostrado anteriormente en la consola.

    \textbf{Sintaxis:}
    \begin{lstlisting}[style=console]
    clear
    \end{lstlisting}

    \subsection{Comando \texttt{exit}}

    Para finalizar la ejecución del programa se utiliza el comando \code{exit}.
    Después de escribirlo y presionar \textit{Enter}, el programa terminará su
    ejecución.

    \textbf{Sintaxis:}
    \begin{lstlisting}[style=console]
    exit
    \end{lstlisting}

    Finalmente, cabe mencionar que el programa incorpora una implementación
    portable de \emph{readline} llamada \emph{(isocline)}, por lo que la
    consola cuenta con historial de comandos navegable mediante las teclas de
    flecha y los atajos de teclado estándar de Bash, como \texttt{Ctrl+R} para
    búsqueda inversa o \texttt{Ctrl+A} y \texttt{Ctrl+E} para moverse al inicio y
    al final de la línea, respectivamente.
\end{document}
